\section{Módulo de Suscripciones}
Este módulo está enfocado en la gestión simplificada de gastos recurrentes, permitiendo al usuario tener un panorama claro de sus compromisos financieros fijos.

\subsection{Indicadores Agregados}
Junto al carril, el módulo expone cuatro cifras: el \textbf{compromiso mensual}, el \textbf{costo anualizado} (doce mensualidades más los cobros anuales), el \textbf{acumulado histórico} efectivamente pagado —con el servicio que más ha costado hasta la fecha— y el número de cobros \textbf{por revisar}.

El acumulado merece mención aparte: un cargo mensual reducido no resulta significativo hasta que se suman los meses transcurridos desde el alta, y es precisamente esa cifra la que permite juzgar si el servicio sigue justificando su costo. Un carril mensual, por su propia naturaleza, no la deja ver.

\subsection{Alta de Servicios}
El registro de una nueva suscripción está diseñado para ser rápido y de baja fricción. El usuario únicamente necesita ingresar dos datos clave para dar de alta el registro:
\begin{itemize}
    \item \textbf{Nombre:} Identificador del servicio (ej. Netflix, Gimnasio, Spotify).
    \item \textbf{Monto:} Valor del cobro recurrente.
    \item \textbf{Ciclo y fecha de cobro:} Periodicidad (mensual o anual) y día —y mes, en el caso anual— en que se efectúa el cargo. Si el día no existe en un mes determinado, el cobro se registra el último día disponible.
    \item \textbf{Medio de pago:} Instrumento al que llega el cargo. Opcional, se declara una vez y lo heredan todos los gastos generados.
\end{itemize}

\subsection{Seguimiento Mensual}
La lógica de control de este módulo se basa en un sistema de verificación periódica. Cada mes, la interfaz presenta el listado de suscripciones activas junto a un \textbf{sistema de \textit{checklist}}. El usuario interactúa marcando un \textit{check} manual en cada servicio una vez que constata que el cobro ha sido procesado o pagado en el mes en curso. Esta mecánica previene el solapamiento de pagos, el olvido de cancelaciones y los cobros sorpresa en la tarjeta.

Marcar un servicio registra automáticamente el gasto correspondiente en el libro mayor. Destildarlo se interpreta como la corrección de un error de registro y, por consiguiente, \textbf{elimina también ese gasto} del historial.

El cierre de mes se ejecuta mediante una acción explícita de \textbf{reseteo}, que desmarca la totalidad de los servicios para comenzar un nuevo ciclo de verificación. Cabe señalar que el estado de pago es un indicador único y no lleva asociado el mes al que corresponde: el sistema no conserva, por tanto, un registro de qué meses fueron pagados más allá de los gastos generados en el libro mayor, y el reseteo depende de que el usuario lo ejecute.